% ID: SEL-D0065-OA-E02
% ESTADO: cotejada. Segunda pasada visual y borrador compilado revisados.
% Fuente: Selectividad/sel_2010_matematicas_aplicadas/mat_apli/mat_apli_sept/mat_apli_sept_ra/EXAMEN-ANDALUCÍA-1.pdf
% SHA-256: 35f05b19d43a5c26840240ff45ede93938ac575e4378437433e833cac887a341
% Página: 1; opción A; ejercicio original 2.

En una empresa han hecho un estudio sobre la rentabilidad de su inversión en publicidad, y han llegado a la conclusión de que el beneficio obtenido, en miles de euros, viene dado por la expresión $B(x)=0.5x^2-4x+6$,

siendo $x$ la inversión en publicidad, en miles de euros, con $x$ en el intervalo $[0,10]$.
\begin{enumerate}[label=\alph*)]
\item \textbf{(1 punto)} \textquestiondown{}Para qué valores de la inversión la empresa tiene pérdidas?
\item \textbf{(1 punto)} \textquestiondown{}Cuánto tiene que invertir la empresa en publicidad para obtener el mayor beneficio posible?
\item \textbf{(0.5 puntos)} \textquestiondown{}Cuál es el beneficio si no se invierte nada en publicidad? \textquestiondown{}Hay algún otro valor de la inversión para el cual se obtiene el mismo beneficio?
\end{enumerate}
